\chapter*{Conclusion}
\addcontentsline{toc}{chapter}{Conclusion}

The book began with a counter that could not go below zero and ended with a search over every small ledger for the questions a summary cannot answer. Between them it built one program in six stages, and the same distinction was drawn at each stage: what the compiler guarantees, what a library guarantees under stated assumptions, and what the program merely arranges and checks by tests. The distinction is the book's main claim about the language. Go gives a programmer types that cannot be forged outside their package, interfaces that name only what a consumer needs, errors that are values to be examined, and goroutines that are cheap enough to use freely, and it leaves unenforced almost everything else: that errors are examined, that a map is not shared, that a function is pure, that a callback does not block. Good Go is mostly a matter of arranging the code so that the unenforced things become visible and testable, and of saying in the documentation which ones remain a matter of discipline.

The same distinction governs the analogies to the author's other work. The staging of inscription, validation, admission, projection, and activation, the separation of routing from adjudication, and the observation that a summary can discard what a later gate requires are structural analogies. Each was implemented in a small program and tested, and none of the formal properties of the specifications they resemble were claimed or proved here. Where a proposition appears in the text, it is a proposition about this code under the stated assumptions, and where the assumptions are not enforced, the text says so. Appendix~\ref{app:guarantees} collects the statements in one table, and its second column is the honest summary of the book: a short list of things the toolchain does for the programmer, and a longer list of things the programmer must keep doing.

Several limits remain open and are named in the text where they arise: the ledger is tamper-evident only relative to a head digest held elsewhere, two processes appending to one file can still fork it, the service speaks plain HTTP, the index file is trusted beyond its position in the ledger, and the manifest identifies a build only as far as the toolchain records it. Each could be closed with more machinery, and each closure would raise a new question of the same shape, which is a fair description of how systems of this kind are built.
